\section{通算企業年金}

\begin{fillblank}[年金1 2021 問題1(6)\textcircled{2} 改題]
  企業年金連合会が支給する通算企業年金について、令和4年5月1日以降に中
途脱退して確定給付企業年金の加入者の資格を喪失した者のうち、移換時年齢45 歳未満の者に
対して適用される予定利率は（　　A　　）％である。
\end{fillblank}

\begin{answerbox}
  \begin{enumerate}[A: ]
    \item 1.25
  \end{enumerate}
\end{answerbox}

\begin{fillblank}[年金1 2019 問題1(4)\textcircled{2}]
  企業年金連合会は2016 年10 月より、存続厚生年金基金、企業年金基金および規約型企業年
金を実施する事業主（以下「年金基金等」という。）が支給する年金または一時金につき一定
額が確保されるよう、年金基金等の拠出金を原資として、年金基金等の年金給付等積立金また
は積立金の額を付加するために（　　A　　）を行っている。
\end{fillblank}

\begin{answerbox}
  \begin{enumerate}[A: ]
    \item 共同運用事業
  \end{enumerate}
\end{answerbox}

\begin{fillblank}[年金1 2018 問題1(5)\textcircled{1} 改題]
企業年金連合会の通算企業年金について、令和4年5月1日以降に 58 歳で受換する中途脱退者に適
用される予定利率は、（　　A　　）％である。
\end{fillblank}

\begin{answerbox}
  \begin{enumerate}[A: ]
    \item 0.75
  \end{enumerate}
\end{answerbox}

\begin{fillblank}[年金1 2018 問題1(5)\textcircled{2}]
企業年金連合会の通算企業年金について、65 歳以降に支給が開始された場合を除き、保証
期間は（　　A　　）とされている。
\end{fillblank}

\begin{answerbox}
  \begin{enumerate}[A: ]
    \item 支給が開始された後、80歳に達するまでの期間
  \end{enumerate}
\end{answerbox}

\begin{fillblank}[年金1 2018 問題1(5)\textcircled{3}]
企業年金連合会が給付の支給に関する権利義務を負っている中途脱退者等が企業型年金加入
者の資格を取得したときは、（　　A　　）を経過する日までに企業年金連合会に申し出ることで、
積立金を移換することができる。
\end{fillblank}

\begin{answerbox}
  \begin{enumerate}[A: ]
    \item 資格を得た日から起算して3ヶ月
  \end{enumerate}
\end{answerbox}

\newpage